Минимизация в каждой ячейке идёт в два этапа: до нормы сил $10^{-3}$
и до $3\cdot10^{-4}$~эВ/\AA. Между этапами напряжения в слоях дальше 5~\AA{}
от поверхности включения меняются не больше чем на «G1_MAX»~МПа.

Ячейка $+\varepsilon$ начинает минимизацию с решения теории упругости,
ячейка $-\varepsilon$ с положений атомов отсчётной ячейки. Поэтому
полусумма откликов на $+\varepsilon$ и $-\varepsilon$ показывает сразу
и нелинейность, и зависимость результата от начального приближения.
В слоях дальше 5~\AA{} от поверхности она не больше «G3_MAX»~МПа, то есть
около 1\% напряжений вдоль поля и поперёк него. В основной статье
используется полуразность откликов.

Атомы матрицы не перескакивают в новые положения. Число атомов, смещение
которых отличается от решения теории упругости больше чем на 0{,}1~\AA,
равно «EV_PLUS» в ячейке $+\varepsilon$ и «EV_MINUS» в ячейке $-\varepsilon$.
Все такие атомы лежат не дальше «EV_MINUS_D»~\AA{} от поверхности
и отклоняются не больше чем на «EV_MINUS_MAX»~\AA; ни один атом не
отклоняется так в обеих ячейках.

Дефекты ищутся методом DXA (алгоритмом выделения дислокаций)
\cite{Stukowski2012DXA}, который находит линии дислокаций по искажению
решётки вокруг них, и анализом ближайших соседей, который определяет тип
кристаллического окружения каждого атома. Поиск проводится в отсчётной
ячейке и в ячейках $\pm\varepsilon$, после каждого этапа минимизации,
с атомами включения и без них. Отрезков дислокаций нет ни в одной ячейке.
Атомы алюминия без окружения ГЦК (гранецентрированной кубической решётки
алюминия) дальше 6~\AA{} от
поверхности во всех трёх ячейках почти одни и те же: «NONFCC_FAR» атомов
в 6--7{,}5~\AA{} от неё, то есть на краю слоя у границы.

Среднее добавленное напряжение по атомам матрицы ($xx$, $yy$, $zz$) равно
$«G4_MD»$~МПа; решение теории упругости для той же ячейки даёт
$«G4_CT»$~МПа. Их отношение, около «G4_RATIO», близко к множителю
«A_FIT», с которым атомистическое поле совпадает с решением теории
упругости (раздел~3.2 основной статьи).

Атомистическое поле слабее решения теории упругости из-за податливого слоя
алюминия у границы включения. Смещения атомов
матрицы, отнесённые к смещениям в решении теории упругости для той же ячейки,
равны по нормали к поверхности включения и вдоль неё «KIN_RAD»\% и «KIN_TAN»\%
в первом слое ($d<2$~\AA), «D46R»\% и «D46T»\% при $d=4$--6~\AA{}
и «D1214R»\% и «D1214T»\% при $d=12$--14~\AA. Деформация каждого атома
найдена по смещениям его соседей в пределах 3{,}2~\AA{} (двенадцать ближайших в ГЦК). При $d>8$~\AA{}
напряжение по вириалу равно кубическим упругим постоянным потенциала,
умноженным на эту деформацию, с наклоном «VIR_FAR» (коэффициент корреляции
не ниже «VIR_CORR»), а при $d<3$~\AA{} наклон равен «VIR_NEAR»: слой у границы
несёт меньшее напряжение при той же деформации. Плотность алюминия при
$r=25{,}0$--$25{,}5$~\AA{} составляет «DENS_FIRST» объёмной. Полость, оставшаяся
для включения после удаления атомов алюминия ближе 2{,}2~\AA{} к нему, больше
включения по объёму на «CAV_PCT»\%. По одному дальнему полю такой податливый слой
нельзя отличить от идеально сцепленной сферы меньшего радиуса: множителю
«A_FIT» соответствует радиус «R_EQ»~\AA. Касательное смещение в первом слое
составляет «KIN_TAN»\% упругого, поэтому при 0~К граница свободно не
проскальзывает.

При деформации включения в 4 и 10 раз больше заданной (0{,}78 и 1{,}94\%,
минимизация до $10^{-3}$~эВ/\AA) дефектов тоже нет: DXA не находит
дислокаций, атомов без окружения ГЦК дальше 6~\AA{} от поверхности не больше,
чем в отсчётной ячейке, а поле, делённое на кратность деформации, совпадает
с полем при $\varepsilon$ (медианное отношение «X4_LIN» и «X10_LIN»). При
$10\varepsilon$ касательное напряжение в лепестках достигает «X10_SHEAR»~МПа.
Расчёт статический и выполнен при 0~К, поэтому термически активированное
зарождение дислокаций он не учитывает. В ячейке нет и исходных дислокаций,
движение которых определяет предел текучести сплава, поэтому отсутствие
дефектов не означает, что сплав у такого включения деформируется только
упруго.
